% Печать notation-glossary (записи: glossary/notation-entries.tex).
\chapter*{Условные обозначения}
\addcontentsline{toc}{chapter}{Условные обозначения}
\label{ch:notation}

% Все записи type=notation, порядок по полю sort (не noidx: иначе пусто до «следующего» прогона).
% В таблице обозначений символы уже с расшифровкой — без pdftooltip (и без \pdfstringdef в math).
\begingroup
\RenewDocumentCommand \gls { m } { \glsentryname {#1} }
\printunsrtglossary[type=notation,style=notationlong,sort=standard,nonumberlist]
\endgroup
